\documentclass[12pt,a4paper]{article}
\usepackage[utf8]{vietnam}
\usepackage{amsmath,amssymb}
\usepackage{tabularx}
\usepackage{tikz}
\usepackage{xcolor}
\usepackage[top=0.6cm, bottom=0.6cm, left=0.6cm, right=0.6cm]{geometry}
\usepackage[dethi]{ex_test}

\SAOPTN{widthSA=2.7,heightSA=0.62}

\setlength{\parskip}{0pt}

% Dòng kẻ chấm mờ đều đặn chuẩn viết tay (khoảng cách 6.5mm rộng rãi vừa vặn)
\newcommand{\dongke}[1]{%
  \begingroup%
  \setlength{\parskip}{0pt}%
  \baselineskip=18.5pt%
  \vspace{1.0mm}%
  \foreach \i in {1,...,#1}{%
    \noindent{\color{gray!60}\dotfill}\par%
  }%
  \vspace{0.6mm}%
  \endgroup%
}

% Định nghĩa ô bài tập tự luận trong bảng
\newcommand{\oBaiTap}[3][6]{%
  \vspace{1.0mm}%
  \noindent\textbf{#2} #3\par%
  \dongke{#1}%
}

\begin{document}
\pagestyle{empty}
% =============================================================================
% ========================= TOÁN 6 - TUẦN 08 - BUỔI 02: MÃ ĐỀ 101 =============
% ========================= MẶT TRƯỚC (TRANG 1/2) ==============================
% =============================================================================

\noindent
\begin{tabularx}{\linewidth}{|X|c|}
\hline
\textbf{PHIẾU KIỂM TRA ĐỊNH KỲ TOÁN 6 - TUẦN 08 - BUỔI 02} & \textbf{Mã đề: 101 - Trang 1/2} \\
\hline
\end{tabularx}

\vspace{1mm}

\noindent
\begin{tabularx}{\linewidth}{|p{3.2cm}|X|}
\hline
\centering\textbf{ĐIỂM SỐ} & \centering\textbf{LỜI NHẬN XÉT CỦA THẦY CÔ} \tabularnewline
\hline
\rule{0pt}{1.8cm} & \tabularnewline
\hline
\end{tabularx}

\vspace{1.2mm}

\noindent\textbf{Câu 1 (0,25 điểm).} Tập hợp $B = \{x \in \mathbb{N}^* \mid x \le 48\}$ có bao nhiêu phần tử? \shortans{48}\par\vspace{0.3mm}
\noindent\textbf{Câu 2 (0,25 điểm).} Số La Mã $\text{XXVI}$ có giá trị trong hệ thập phân là bao nhiêu? \shortans{26}\par\vspace{0.3mm}
\noindent\textbf{Câu 3 (0,25 điểm).} Viết gọn kết quả của phép chia $5^8 : 5^3$ dưới dạng một lũy thừa. \shortans{5^5}\par\vspace{0.3mm}
\noindent\textbf{Câu 4 (0,25 điểm).} Giá trị của biểu thức $40 + 20 : 2^2$ bằng bao nhiêu? \shortans{45}\par\vspace{0.3mm}
\noindent\textbf{Câu 5 (0,25 điểm).} Chữ số tận cùng của số tự nhiên chia hết cho cả $2$ và $5$ là: \shortans{0}\par\vspace{0.3mm}
\noindent\textbf{Câu 6 (0,25 điểm).} Tìm chữ số $x$ nhỏ nhất khác $0$ để số $\overline{5x2}$ chia hết cho $3$. \shortans{2}\par\vspace{0.3mm}
\noindent\textbf{Câu 7 (0,25 điểm).} Phân tích số $24$ ra thừa số nguyên tố có kết quả là gì? \shortans{2^3 \cdot 3}\par\vspace{0.3mm}
\noindent\textbf{Câu 8 (0,25 điểm).} Hai đường chéo của hình bình hành cắt nhau tại: \shortans{\text{Trung điểm}}\par\vspace{0.3mm}
\noindent\textbf{Câu 9 (0,25 điểm).} Một hình chữ nhật có chiều dài $14\text{ m}$, chiều rộng $6\text{ m}$. Chu vi là: \shortans{40\text{ m}}\par\vspace{0.3mm}
\noindent\textbf{Câu 10 (0,25 điểm).} Hình bình hành có độ dài đáy $15\text{ cm}$, chiều cao $8\text{ cm}$. Diện tích là: \shortans{120\text{ cm}^2}\par\vspace{0.3mm}
\vspace{0.4mm}

\noindent\textbf{Câu 11 (3,0 điểm). Thực hiện phép tính (hợp lý nếu có thể):}
\vspace{0.6mm}

\noindent
\begin{tabularx}{\linewidth}{|X|X|}
\hline
\oBaiTap[6]{a)}{185 - 3^2 \cdot 6 + 2^4 \cdot 3} & 
\oBaiTap[6]{b)}{125 + 74 + 175 + 26} \\
\hline
\oBaiTap[6]{c)}{37 \cdot 64 + 37 \cdot 36 + 300} & 
\oBaiTap[6]{d)}{240 : (35 - 11) + 2^3 \cdot 5} \\
\hline
\oBaiTap[6]{e)}{160 : [52 - (15 - 9)^2]} & 
\oBaiTap[6]{f)}{500 : \{ 130 - [53 + (17 - 14)^3] \}} \\
\hline
\end{tabularx}

\newpage
% =============================================================================
% ========================= TOÁN 6 - TUẦN 08 - BUỔI 02: MÃ ĐỀ 101 =============
% ========================= MẶT SAU (TRANG 2/2) ================================
% =============================================================================

\noindent
\begin{tabularx}{\linewidth}{|X|c|}
\hline
\textbf{TOÁN 6 - TUẦN 08 - BUỔI 02} & \textbf{Mã đề: 101 - Trang 2/2} \\
\hline
\end{tabularx}

\vspace{1.5mm}

\noindent\textbf{Câu 12 (2,0 điểm). Tìm số tự nhiên $x$, biết:}
\vspace{0.8mm}

\noindent
\begin{tabularx}{\linewidth}{|X|X|}
\hline
\oBaiTap[5]{a)}{(x - 12) : 4 = 15} & 
\oBaiTap[5]{b)}{120 - [45 - (x - 3)] = 85} \\
\hline
\oBaiTap[5]{c)}{(x - 2)^3 = 27} & 
\oBaiTap[5]{d)}{$x \ \vdots \ 12$, $x \ \vdots \ 18$ và $50 < x < 100$.} \\
\hline
\end{tabularx}

\vspace{1.5mm}

\noindent
\begin{tabularx}{\linewidth}{|X|}
\hline
\vspace{0.5mm}
\textbf{Câu 13 (1,5 điểm).} Chi đội $6\text{A}$ có $24$ bạn nam và $36$ bạn nữ. Cô giáo chủ nhiệm muốn chia đều số nam và số nữ vào các tổ học tập sao cho số nam và số nữ ở mỗi tổ là như nhau.\par
\textbf{a) (0,75 điểm)} Cô giáo có thể chia được nhiều nhất thành bao nhiêu tổ?\par
\textbf{b) (0,75 điểm)} Khi đó, ở mỗi tổ có bao nhiêu bạn nam và bao nhiêu bạn nữ?\par
\vspace{0.4mm}
\noindent\textit{Bài giải:}\dongke{12}
\\ \hline
\end{tabularx}

\vspace{1.5mm}

\noindent
\begin{tabularx}{\linewidth}{|X|}
\hline
\vspace{0.5mm}
\textbf{Câu 14 (1,0 điểm).} Cho biểu thức $A = 2 + 2^2 + 2^3 + \dots + 2^{60}$. Chứng minh rằng $A$ chia hết cho $7$.\par
\vspace{0.4mm}
\noindent\textit{Bài giải:}\dongke{9}
\\ \hline
\end{tabularx}

\newpage

% =============================================================================
% ========================= TOÁN 6 - TUẦN 08 - BUỔI 02: MÃ ĐỀ 102 =============
% ========================= MẶT TRƯỚC (TRANG 1/2) ==============================
% =============================================================================

\noindent
\begin{tabularx}{\linewidth}{|X|c|}
\hline
\textbf{PHIẾU KIỂM TRA ĐỊNH KỲ TOÁN 6 - TUẦN 08 - BUỔI 02} & \textbf{Mã đề: 102 - Trang 1/2} \\
\hline
\end{tabularx}

\vspace{1mm}

\noindent
\begin{tabularx}{\linewidth}{|p{3.2cm}|X|}
\hline
\centering\textbf{ĐIỂM SỐ} & \centering\textbf{LỜI NHẬN XÉT CỦA THẦY CÔ} \tabularnewline
\hline
\rule{0pt}{1.8cm} & \tabularnewline
\hline
\end{tabularx}

\vspace{1.2mm}

\noindent\textbf{Câu 1 (0,25 điểm).} Tập hợp $B = \{x \in \mathbb{N}^* \mid x \le 56\}$ có bao nhiêu phần tử? \shortans{56}\par\vspace{0.3mm}
\noindent\textbf{Câu 2 (0,25 điểm).} Số La Mã $\text{XXIII}$ có giá trị trong hệ thập phân là bao nhiêu? \shortans{23}\par\vspace{0.3mm}
\noindent\textbf{Câu 3 (0,25 điểm).} Viết gọn kết quả của phép chia $6^9 : 6^4$ dưới dạng một lũy thừa. \shortans{6^5}\par\vspace{0.3mm}
\noindent\textbf{Câu 4 (0,25 điểm).} Giá trị của biểu thức $50 + 36 : 3^2$ bằng bao nhiêu? \shortans{54}\par\vspace{0.3mm}
\noindent\textbf{Câu 5 (0,25 điểm).} Số dư khi chia một số có chữ số tận cùng là $3$ cho $5$ là: \shortans{3}\par\vspace{0.3mm}
\noindent\textbf{Câu 6 (0,25 điểm).} Tìm chữ số $x$ nhỏ nhất khác $0$ để số $\overline{4x5}$ chia hết cho $3$. \shortans{3}\par\vspace{0.3mm}
\noindent\textbf{Câu 7 (0,25 điểm).} Phân tích số $36$ ra thừa số nguyên tố có kết quả là gì? \shortans{2^2 \cdot 3^2}\par\vspace{0.3mm}
\noindent\textbf{Câu 8 (0,25 điểm).} Hai cạnh đối diện của hình bình hành có tính chất: \shortans{\text{Song song, bằng nhau}}\par\vspace{0.3mm}
\noindent\textbf{Câu 9 (0,25 điểm).} Một hình chữ nhật có chiều dài $15\text{ m}$, chiều rộng $7\text{ m}$. Chu vi là: \shortans{44\text{ m}}\par\vspace{0.3mm}
\noindent\textbf{Câu 10 (0,25 điểm).} Hình bình hành có độ dài đáy $16\text{ cm}$, chiều cao $7\text{ cm}$. Diện tích là: \shortans{112\text{ cm}^2}\par\vspace{0.3mm}
\vspace{0.4mm}

\noindent\textbf{Câu 11 (3,0 điểm). Thực hiện phép tính (hợp lý nếu có thể):}
\vspace{0.6mm}

\noindent
\begin{tabularx}{\linewidth}{|X|X|}
\hline
\oBaiTap[6]{a)}{195 - 4^2 \cdot 4 + 3^3 \cdot 2} & 
\oBaiTap[6]{b)}{135 + 68 + 165 + 32} \\
\hline
\oBaiTap[6]{c)}{43 \cdot 78 + 43 \cdot 22 + 700} & 
\oBaiTap[6]{d)}{260 : (38 - 12) + 3^2 \cdot 6} \\
\hline
\oBaiTap[6]{e)}{180 : [67 - (17 - 10)^2]} & 
\oBaiTap[6]{f)}{600 : \{ 140 - [53 + (18 - 15)^3] \}} \\
\hline
\end{tabularx}

\newpage
% =============================================================================
% ========================= TOÁN 6 - TUẦN 08 - BUỔI 02: MÃ ĐỀ 102 =============
% ========================= MẶT SAU (TRANG 2/2) ================================
% =============================================================================

\noindent
\begin{tabularx}{\linewidth}{|X|c|}
\hline
\textbf{TOÁN 6 - TUẦN 08 - BUỔI 02} & \textbf{Mã đề: 102 - Trang 2/2} \\
\hline
\end{tabularx}

\vspace{1.5mm}

\noindent\textbf{Câu 12 (2,0 điểm). Tìm số tự nhiên $x$, biết:}
\vspace{0.8mm}

\noindent
\begin{tabularx}{\linewidth}{|X|X|}
\hline
\oBaiTap[5]{a)}{(x - 15) : 5 = 14} & 
\oBaiTap[5]{b)}{130 - [55 - (x - 4)] = 90} \\
\hline
\oBaiTap[5]{c)}{(x - 3)^3 = 64} & 
\oBaiTap[5]{d)}{$x \ \vdots \ 15$, $x \ \vdots \ 20$ và $80 < x < 150$.} \\
\hline
\end{tabularx}

\vspace{1.5mm}

\noindent
\begin{tabularx}{\linewidth}{|X|}
\hline
\vspace{0.5mm}
\textbf{Câu 13 (1,5 điểm).} Chi đội $6\text{B}$ có $28$ bạn nam và $42$ bạn nữ. Cô giáo chủ nhiệm muốn chia đều số nam và số nữ vào các tổ học tập sao cho số tổ là nhiều nhất.\par
\textbf{a) (0,75 điểm)} Cô giáo có thể chia được nhiều nhất thành bao nhiêu tổ?\par
\textbf{b) (0,75 điểm)} Khi đó, ở mỗi tổ có bao nhiêu bạn nam và bao nhiêu bạn nữ?\par
\vspace{0.4mm}
\noindent\textit{Bài giải:}\dongke{12}
\\ \hline
\end{tabularx}

\vspace{1.5mm}

\noindent
\begin{tabularx}{\linewidth}{|X|}
\hline
\vspace{0.5mm}
\textbf{Câu 14 (1,0 điểm).} Cho biểu thức $B = 3 + 3^2 + 3^3 + \dots + 3^{60}$. Chứng minh rằng $B$ chia hết cho $13$.\par
\vspace{0.4mm}
\noindent\textit{Bài giải:}\dongke{9}
\\ \hline
\end{tabularx}

\newpage

% =============================================================================
% ========================= TOÁN 6 - TUẦN 08 - BUỔI 02: MÃ ĐỀ 103 =============
% ========================= MẶT TRƯỚC (TRANG 1/2) ==============================
% =============================================================================

\noindent
\begin{tabularx}{\linewidth}{|X|c|}
\hline
\textbf{PHIẾU KIỂM TRA ĐỊNH KỲ TOÁN 6 - TUẦN 08 - BUỔI 02} & \textbf{Mã đề: 103 - Trang 1/2} \\
\hline
\end{tabularx}

\vspace{1mm}

\noindent
\begin{tabularx}{\linewidth}{|p{3.2cm}|X|}
\hline
\centering\textbf{ĐIỂM SỐ} & \centering\textbf{LỜI NHẬN XÉT CỦA THẦY CÔ} \tabularnewline
\hline
\rule{0pt}{1.8cm} & \tabularnewline
\hline
\end{tabularx}

\vspace{1.2mm}

\noindent\textbf{Câu 1 (0,25 điểm).} Tập hợp $B = \{x \in \mathbb{N}^* \mid x \le 64\}$ có bao nhiêu phần tử? \shortans{64}\par\vspace{0.3mm}
\noindent\textbf{Câu 2 (0,25 điểm).} Số La Mã $\text{XXV}$ có giá trị trong hệ thập phân là bao nhiêu? \shortans{25}\par\vspace{0.3mm}
\noindent\textbf{Câu 3 (0,25 điểm).} Viết gọn kết quả của phép chia $7^7 : 7^2$ dưới dạng một lũy thừa. \shortans{7^5}\par\vspace{0.3mm}
\noindent\textbf{Câu 4 (0,25 điểm).} Giá trị của biểu thức $60 + 48 : 2^4$ bằng bao nhiêu? \shortans{63}\par\vspace{0.3mm}
\noindent\textbf{Câu 5 (0,25 điểm).} Số dư khi chia một số có chữ số tận cùng là $7$ cho $5$ là: \shortans{2}\par\vspace{0.3mm}
\noindent\textbf{Câu 6 (0,25 điểm).} Tìm chữ số $x$ nhỏ nhất khác $0$ để số $\overline{6x1}$ chia hết cho $3$. \shortans{2}\par\vspace{0.3mm}
\noindent\textbf{Câu 7 (0,25 điểm).} Phân tích số $40$ ra thừa số nguyên tố có kết quả là gì? \shortans{2^3 \cdot 5}\par\vspace{0.3mm}
\noindent\textbf{Câu 8 (0,25 điểm).} Hai góc đối diện của một hình bình hành có số đo: \shortans{\text{Bằng nhau}}\par\vspace{0.3mm}
\noindent\textbf{Câu 9 (0,25 điểm).} Một hình chữ nhật có chiều dài $16\text{ m}$, chiều rộng $8\text{ m}$. Chu vi là: \shortans{48\text{ m}}\par\vspace{0.3mm}
\noindent\textbf{Câu 10 (0,25 điểm).} Hình bình hành có độ dài đáy $18\text{ cm}$, chiều cao $6\text{ cm}$. Diện tích là: \shortans{108\text{ cm}^2}\par\vspace{0.3mm}
\vspace{0.4mm}

\noindent\textbf{Câu 11 (3,0 điểm). Thực hiện phép tính (hợp lý nếu có thể):}
\vspace{0.6mm}

\noindent
\begin{tabularx}{\linewidth}{|X|X|}
\hline
\oBaiTap[6]{a)}{175 - 2^5 \cdot 3 + 5^2 \cdot 3} & 
\oBaiTap[6]{b)}{145 + 57 + 155 + 43} \\
\hline
\oBaiTap[6]{c)}{52 \cdot 83 + 52 \cdot 17 + 800} & 
\oBaiTap[6]{d)}{280 : (42 - 14) + 4^2 \cdot 3} \\
\hline
\oBaiTap[6]{e)}{200 : [84 - (19 - 11)^2]} & 
\oBaiTap[6]{f)}{700 : \{ 150 - [53 + (19 - 16)^3] \}} \\
\hline
\end{tabularx}

\newpage
% =============================================================================
% ========================= TOÁN 6 - TUẦN 08 - BUỔI 02: MÃ ĐỀ 103 =============
% ========================= MẶT SAU (TRANG 2/2) ================================
% =============================================================================

\noindent
\begin{tabularx}{\linewidth}{|X|c|}
\hline
\textbf{TOÁN 6 - TUẦN 08 - BUỔI 02} & \textbf{Mã đề: 103 - Trang 2/2} \\
\hline
\end{tabularx}

\vspace{1.5mm}

\noindent\textbf{Câu 12 (2,0 điểm). Tìm số tự nhiên $x$, biết:}
\vspace{0.8mm}

\noindent
\begin{tabularx}{\linewidth}{|X|X|}
\hline
\oBaiTap[5]{a)}{(x - 14) : 6 = 12} & 
\oBaiTap[5]{b)}{140 - [65 - (x - 5)] = 95} \\
\hline
\oBaiTap[5]{c)}{(x - 1)^3 = 125} & 
\oBaiTap[5]{d)}{$x \ \vdots \ 14$, $x \ \vdots \ 21$ và $60 < x < 100$.} \\
\hline
\end{tabularx}

\vspace{1.5mm}

\noindent
\begin{tabularx}{\linewidth}{|X|}
\hline
\vspace{0.5mm}
\textbf{Câu 13 (1,5 điểm).} Chi đội $6\text{C}$ có $32$ bạn nam và $48$ bạn nữ. Cô giáo chủ nhiệm muốn chia đều số nam và số nữ vào các nhóm sao cho số nhóm là nhiều nhất.\par
\textbf{a) (0,75 điểm)} Cô giáo có thể chia được nhiều nhất thành bao nhiêu nhóm?\par
\textbf{b) (0,75 điểm)} Khi đó, ở mỗi nhóm có bao nhiêu bạn nam và bao nhiêu bạn nữ?\par
\vspace{0.4mm}
\noindent\textit{Bài giải:}\dongke{12}
\\ \hline
\end{tabularx}

\vspace{1.5mm}

\noindent
\begin{tabularx}{\linewidth}{|X|}
\hline
\vspace{0.5mm}
\textbf{Câu 14 (1,0 điểm).} Cho biểu thức $C = 4 + 4^2 + 4^3 + \dots + 4^{60}$. Chứng minh rằng $C$ chia hết cho $21$.\par
\vspace{0.4mm}
\noindent\textit{Bài giải:}\dongke{9}
\\ \hline
\end{tabularx}

\newpage

% =============================================================================
% ========================= TOÁN 6 - TUẦN 08 - BUỔI 02: MÃ ĐỀ 104 =============
% ========================= MẶT TRƯỚC (TRANG 1/2) ==============================
% =============================================================================

\noindent
\begin{tabularx}{\linewidth}{|X|c|}
\hline
\textbf{PHIẾU KIỂM TRA ĐỊNH KỲ TOÁN 6 - TUẦN 08 - BUỔI 02} & \textbf{Mã đề: 104 - Trang 1/2} \\
\hline
\end{tabularx}

\vspace{1mm}

\noindent
\begin{tabularx}{\linewidth}{|p{3.2cm}|X|}
\hline
\centering\textbf{ĐIỂM SỐ} & \centering\textbf{LỜI NHẬN XÉT CỦA THẦY CÔ} \tabularnewline
\hline
\rule{0pt}{1.8cm} & \tabularnewline
\hline
\end{tabularx}

\vspace{1.2mm}

\noindent\textbf{Câu 1 (0,25 điểm).} Tập hợp $B = \{x \in \mathbb{N}^* \mid x \le 75\}$ có bao nhiêu phần tử? \shortans{75}\par\vspace{0.3mm}
\noindent\textbf{Câu 2 (0,25 điểm).} Số La Mã $\text{XXII}$ có giá trị trong hệ thập phân là bao nhiêu? \shortans{22}\par\vspace{0.3mm}
\noindent\textbf{Câu 3 (0,25 điểm).} Viết gọn kết quả của phép chia $8^8 : 8^3$ dưới dạng một lũy thừa. \shortans{8^5}\par\vspace{0.3mm}
\noindent\textbf{Câu 4 (0,25 điểm).} Giá trị của biểu thức $70 + 75 : 5^2$ bằng bao nhiêu? \shortans{73}\par\vspace{0.3mm}
\noindent\textbf{Câu 5 (0,25 điểm).} Số dư khi chia một số có chữ số tận cùng là $8$ cho $5$ là: \shortans{3}\par\vspace{0.3mm}
\noindent\textbf{Câu 6 (0,25 điểm).} Tìm chữ số $x$ nhỏ nhất khác $0$ để số $\overline{7x3}$ chia hết cho $3$. \shortans{2}\par\vspace{0.3mm}
\noindent\textbf{Câu 7 (0,25 điểm).} Phân tích số $48$ ra thừa số nguyên tố có kết quả là gì? \shortans{2^4 \cdot 3}\par\vspace{0.3mm}
\noindent\textbf{Câu 8 (0,25 điểm).} Hai đường chéo của một hình thoi có tính chất: \shortans{\text{Vuông góc}}\par\vspace{0.3mm}
\noindent\textbf{Câu 9 (0,25 điểm).} Một hình chữ nhật có chiều dài $18\text{ m}$, chiều rộng $9\text{ m}$. Chu vi là: \shortans{54\text{ m}}\par\vspace{0.3mm}
\noindent\textbf{Câu 10 (0,25 điểm).} Hình bình hành có độ dài đáy $20\text{ cm}$, chiều cao $8\text{ cm}$. Diện tích là: \shortans{160\text{ cm}^2}\par\vspace{0.3mm}
\vspace{0.4mm}

\noindent\textbf{Câu 11 (3,0 điểm). Thực hiện phép tính (hợp lý nếu có thể):}
\vspace{0.6mm}

\noindent
\begin{tabularx}{\linewidth}{|X|X|}
\hline
\oBaiTap[6]{a)}{165 - 3^3 \cdot 3 + 4^3 \cdot 1} & 
\oBaiTap[6]{b)}{115 + 86 + 185 + 14} \\
\hline
\oBaiTap[6]{c)}{61 \cdot 74 + 61 \cdot 26 + 900} & 
\oBaiTap[6]{d)}{300 : (45 - 15) + 5^2 \cdot 2} \\
\hline
\oBaiTap[6]{e)}{220 : [103 - (21 - 12)^2]} & 
\oBaiTap[6]{f)}{800 : \{ 160 - [53 + (20 - 17)^3] \}} \\
\hline
\end{tabularx}

\newpage
% =============================================================================
% ========================= TOÁN 6 - TUẦN 08 - BUỔI 02: MÃ ĐỀ 104 =============
% ========================= MẶT SAU (TRANG 2/2) ================================
% =============================================================================

\noindent
\begin{tabularx}{\linewidth}{|X|c|}
\hline
\textbf{TOÁN 6 - TUẦN 08 - BUỔI 02} & \textbf{Mã đề: 104 - Trang 2/2} \\
\hline
\end{tabularx}

\vspace{1.5mm}

\noindent\textbf{Câu 12 (2,0 điểm). Tìm số tự nhiên $x$, biết:}
\vspace{0.8mm}

\noindent
\begin{tabularx}{\linewidth}{|X|X|}
\hline
\oBaiTap[5]{a)}{(x - 16) : 7 = 11} & 
\oBaiTap[5]{b)}{150 - [75 - (x - 6)] = 100} \\
\hline
\oBaiTap[5]{c)}{(x - 4)^3 = 8} & 
\oBaiTap[5]{d)}{$x \ \vdots \ 16$, $x \ \vdots \ 24$ và $70 < x < 120$.} \\
\hline
\end{tabularx}

\vspace{1.5mm}

\noindent
\begin{tabularx}{\linewidth}{|X|}
\hline
\vspace{0.5mm}
\textbf{Câu 13 (1,5 điểm).} Chi đội $6\text{D}$ có $30$ bạn nam và $45$ bạn nữ. Cô giáo chủ nhiệm muốn chia đều số nam và số nữ vào các tổ học tập sao cho số tổ là nhiều nhất.\par
\textbf{a) (0,75 điểm)} Cô giáo có thể chia được nhiều nhất thành bao nhiêu tổ?\par
\textbf{b) (0,75 điểm)} Khi đó, ở mỗi tổ có bao nhiêu bạn nam và bao nhiêu bạn nữ?\par
\vspace{0.4mm}
\noindent\textit{Bài giải:}\dongke{12}
\\ \hline
\end{tabularx}

\vspace{1.5mm}

\noindent
\begin{tabularx}{\linewidth}{|X|}
\hline
\vspace{0.5mm}
\textbf{Câu 14 (1,0 điểm).} Cho biểu thức $D = 5 + 5^2 + 5^3 + \dots + 5^{60}$. Chứng minh rằng $D$ chia hết cho $31$.\par
\vspace{0.4mm}
\noindent\textit{Bài giải:}\dongke{9}
\\ \hline
\end{tabularx}

\newpage

\end{document}